\documentclass[11pt]{article}
\usepackage[margin=1in]{geometry}
\usepackage{amsmath,amssymb,amsthm}
\usepackage{algorithm}
\usepackage{algpseudocode}
\usepackage{booktabs}
\usepackage{hyperref}

\newtheorem{definition}{Definition}
\newtheorem{theorem}{Theorem}
\newtheorem{lemma}{Lemma}
\newtheorem{example}{Example}

\title{Lexicographic Reverse Search:\\Extreme-Point Enumeration of Local Credal Sets}
\author{IBM Research}
\date{}

\begin{document}
\maketitle

\begin{abstract}
The credal-network family of inference engines in this library (CredalVE, IntervalBP,
CredalCTE, ApproxLP, CredalIJGP) all consume the \emph{extreme points} (vertices) of the
local credal sets rather than the interval bounds directly. Converting a set of interval
bounds $[\ell_i, u_i]$ on a conditional distribution into the vertex representation of the
corresponding polytope is a \emph{vertex-enumeration} problem, solved here by pyAgrum's
\texttt{intervalToCredal()} which runs the \textbf{Lexicographic Reverse Search (LRS)}
algorithm of Avis and Fukuda. This note describes the geometry of the local credal sets,
the LRS algorithm and its output-sensitive guarantees, works a complete example on the
compiled credal network of \texttt{alarm.lcn}, and analyses why the enumeration --- cheap
for narrow families --- becomes the dominant cost for wide or high-cardinality families.
\end{abstract}

\section{From interval bounds to vertices}

\subsection{The local credal set of a family}

After an LCN is compiled into a credal network (see \texttt{credal\_network.py}), each node
$X$ with cardinality $d = |\mathcal{D}_X|$ and each configuration $\pi$ of its parents
carries a \textbf{local credal set}: a convex set of conditional distributions
$P(X \mid \pi)$. In the interval representation produced by the compiler, this set is
specified by lower/upper bounds $\ell_i \le P(X{=}i \mid \pi) \le u_i$ for each state
$i \in \{1, \dots, d\}$. Geometrically it is the \emph{interval-constrained probability
simplex}
\begin{equation}
  K(X \mid \pi) \;=\; \Bigl\{\, p \in \mathbb{R}^d \;:\;
    \textstyle\sum_{i=1}^d p_i = 1,\quad
    \ell_i \le p_i \le u_i \ \ (1 \le i \le d) \,\Bigr\}.
  \label{eq:local}
\end{equation}
This is a bounded polytope living in the $(d{-}1)$-dimensional affine hyperplane
$\sum_i p_i = 1$, cut out by the $2d$ interval facets. Downstream engines need its
\textbf{V-representation}: the finite set $\mathrm{ext}\,K(X\mid\pi)$ of vertices, each of
which is a concrete probability distribution over $X$.

\subsection{H- and V-representations}

A convex polytope admits two dual descriptions:
\begin{itemize}
  \item the \textbf{H-representation}: an intersection of half-spaces,
    $\{x : Ax \le b\}$ (here, the $2d$ interval inequalities plus the normalization
    equality) --- this is what the compiler stores;
  \item the \textbf{V-representation}: the convex hull of its vertices,
    $\mathrm{conv}\{v_1, \dots, v_V\}$ --- this is what the engines consume.
\end{itemize}
The \textbf{vertex-enumeration problem} is: given the H-representation, output every
vertex. This is exactly what \texttt{gum.CredalNet.intervalToCredal()} does, once per
$(X, \pi)$ pair, in \texttt{CredalNetworkVertices.\_build}:
\begin{verbatim}
    self.credal_net = gum.CredalNet(self.bn_min, self.bn_max)
    self.credal_net.intervalToCredal()   # LRS vertex enumeration
\end{verbatim}
Note that the enumeration is run \emph{per local credal set}, not on any global joint
polytope: there is one polytope of the form~\eqref{eq:local} for every parent
configuration of every node.

\section{The LRS algorithm}

\subsection{Reverse search}

Vertex enumeration can be done by pivoting (Avis--Fukuda's \emph{lrs}) or by the double
description method (Motzkin's \emph{cdd}). pyAgrum uses \textbf{Lexicographic Reverse
Search (LRS)}~\cite{avisfukuda}, a pivoting method with a decisive advantage: it uses
memory independent of the output size and never stores the vertices it has already found.

Reverse search is a general technique for enumerating the nodes of an implicitly defined
graph. For vertex enumeration the graph is the \textbf{skeleton} of the polytope --- its
vertices connected by edges. The ingredients are:
\begin{enumerate}
  \item A \emph{local search / pivot rule} $f$ that, from any vertex, deterministically
    moves to a unique ``parent'' vertex. Under Bland's/the lexicographic rule this is
    the reverse of the simplex pivot, and following $f$ repeatedly from any vertex always
    terminates at the same unique \emph{root} (the lexicographically-minimal optimal
    vertex of a chosen objective).
  \item Because $f$ is deterministic and acyclic, the ``$f$ leads to parent'' relation
    induces a \textbf{spanning tree} of the skeleton rooted at that vertex.
  \item Reverse search enumerates the tree by DFS \emph{from the root outward}: at each
    vertex it tries every adjacent edge (every reverse pivot); a neighbour is accepted
    as a child iff applying $f$ to it returns to the current vertex.
\end{enumerate}
The elegance is that the tree is never materialised. LRS holds only the current vertex
and the current pivot indices --- $O(d\cdot m)$ working memory --- and reconstructs its
path by pivoting, so it can enumerate astronomically many vertices in bounded space.
The ``lexicographic'' qualifier is the perturbation rule that breaks ties consistently
and thereby handles \textbf{degenerate} polytopes (more than $d$ facets meeting at a
vertex) --- which the interval simplices~\eqref{eq:local} routinely are, since a vertex
is pinned by choosing which of the $2d$ bounds are tight.

\begin{algorithm}[t]
\caption{Reverse search vertex enumeration (schematic)}
\label{alg:lrs}
\begin{algorithmic}[1]
\Procedure{LRS}{$A, b$}\Comment{polytope $\{x : Ax \le b\}$}
  \State $v \gets$ root vertex (lex-min optimum of a generic objective $c^\top x$)
  \State $j \gets 0$ \Comment{index of the reverse pivot being tried}
  \Repeat
    \State \textbf{output} $v$ when first visited
    \While{$j <$ number of candidate pivots at $v$}
      \State $j \gets j+1$
      \State $w \gets \textsc{ReversePivot}(v, j)$
      \If{$w$ is a vertex \textbf{and} $f(w) = v$} \Comment{$w$ is a child}
        \State $v \gets w$; \ $j \gets 0$ \Comment{descend}
        \State \textbf{break}
      \EndIf
    \EndWhile
    \If{no child was found}\Comment{backtrack via the deterministic pivot $f$}
      \State $(v, j) \gets f(v)$ and restore the pivot index at the parent
    \EndIf
  \Until{back at the root with all pivots exhausted}
\EndProcedure
\end{algorithmic}
\end{algorithm}

\subsection{Output-sensitive complexity}

Let a single local polytope~\eqref{eq:local} have
\begin{itemize}
  \item dimension $d' = d-1$ (states $d = |\mathcal{D}_X|$, minus the normalization);
  \item $m \approx 2d$ facets (the interval bounds);
  \item $V$ output vertices.
\end{itemize}
LRS enumerates all $V$ vertices in
\begin{equation}
  O\bigl(d' \cdot m \cdot V\bigr) \ \text{time}, \qquad O\bigl(d' \cdot m\bigr)\ \text{space}.
  \label{eq:lrs-cost}
\end{equation}
That is, the cost per vertex is \emph{polynomial} ($O(d' m)$, one pivot table update) and
the algorithm is \textbf{output-sensitive}: its running time is proportional to the size
$V$ of the answer it produces, with polynomial overhead and no dependence of the
\emph{memory} on $V$. This is optimal in the sense that any algorithm must at least write
its $V$ outputs; LRS adds only a polynomial factor and constant space. (For non-degenerate
input LRS is provably polynomial-total-time; the lexicographic rule extends this to the
degenerate case at the cost of processing some ``dictionaries'' that do not emit a new
vertex, which for the highly-degenerate interval simplices is a real but bounded overhead.)

\section{Worked example: the credal network of \texttt{alarm.lcn}}

\subsection{The model}

The LCN \texttt{examples/alarm.lcn} is
\begin{align*}
  s_1&: 0.1 \le P(B) \le 0.2, & s_2&: 0.05 \le P(E) \le 0.1,\\
  s_3&: 0.6 \le P(A \mid B \vee E) \le 0.9, &
  s_4&: 0.5 \le P(\neg(C \oplus D) \mid A) \le 0.7,\\
  s_5&: 0.05 \le P(A) \le 0.1.
\end{align*}
Compiling it (\texttt{CredalNetwork.from\_lcn}) yields a chain-graph credal network over
\textbf{four nodes} $\{A, B, E, C\text{-}D\}$, where $C$ and $D$ are fused into a single
compound node $C\text{-}D$ of cardinality $4$ because the sentence $s_4$ couples them.
The directed structure is $B \to A$, $E \to A$, and $A \to C\text{-}D$. The compiled
interval local credal sets (one row per family interpretation) are:

\begin{center}
\small
\begin{tabular}{llcc}
\toprule
node & (states, parents) config & $\ell$ & $u$ \\
\midrule
$B$ & $B{=}0$ & 0.8000 & 0.9000 \\
    & $B{=}1$ & 0.1000 & 0.2000 \\
\midrule
$E$ & $E{=}0$ & 0.9000 & 0.9500 \\
    & $E{=}1$ & 0.0500 & 0.1000 \\
\midrule
$A$ & $A{=}0 \mid B{=}0,E{=}0$ & 0.9556 & 1.0000 \\
    & $A{=}1 \mid B{=}0,E{=}0$ & 0.0000 & 0.0444 \\
    & $\;\cdots\mid$ (other $B,E$ configs) & 0.0000 & 1.0000 \\
\midrule
$C\text{-}D$ & $CD \mid A{=}0$ & \multicolumn{2}{c}{$u = (1,0.7,1,0.5)$, $\ell = 0$}\\
             & $CD \mid A{=}1$ & \multicolumn{2}{c}{$u = (1,0.5,1,0.7)$, $\ell = 0$}\\
\bottomrule
\end{tabular}
\end{center}

\subsection{Enumerating each local credal set}

LRS is run once per (node, parent-config) pair. The results
(\texttt{cnv.extreme\_points}) are:

\paragraph{Binary root nodes ($d=2$, $d'=1$).} The polytope is a $1$-dimensional
segment; it has exactly $2$ vertices, its endpoints.
\begin{align*}
  K(B) &= \mathrm{conv}\{(0.8,\,0.2),\ (0.9,\,0.1)\},\\
  K(E) &= \mathrm{conv}\{(0.9,\,0.1),\ (0.95,\,0.05)\}.
\end{align*}

\paragraph{The binary child $A$ ($d=2$).} Each parent configuration is again a segment
with $2$ vertices. For $B{=}0,E{=}0$ the interval $P(A{=}0) \in [0.9556, 1]$ gives
\[
  K(A \mid B{=}0,E{=}0) = \mathrm{conv}\{(0.9556,\,0.0444),\ (1.0,\,0.0)\},
\]
while the three configurations touching $B\vee E$ collapse to the whole edge
$\mathrm{conv}\{(0,1),(1,0)\}$ (the $[0,1]$ interval is vacuous there).

\paragraph{The compound node $C\text{-}D$ ($d=4$, $d'=3$).} This is where enumeration bites.
The local set is the $3$-dimensional interval simplex~\eqref{eq:local} with $m=8$ facets.
\begin{itemize}
  \item Under $A{=}0$ the upper bounds are $(1, 0.7, 1, 0.5)$; LRS returns the
    \textbf{4} unit vertices $\{(0,0,0,1),(0,0,1,0),(0,1,0,0),(1,0,0,0)\}$ (the
    upper bounds are loose enough that only the simplex corners survive).
  \item Under $A{=}1$ the upper bounds are $(1, 0.5, 1, 0.7)$; these cut the simplex
    into a polytope with \textbf{11} vertices, e.g.
    \[
      (0.5,0.5,0,0),\ (0.7,0,0.3,0),\ (0.3,0,0,0.7),\ (0,0.5,0.5,0),\
      (0,0,0.3,0.7),\ \dots
    \]
    each corresponding to a distinct choice of which interval bounds are active.
\end{itemize}

\subsection{Why $C\text{-}D$ is the expensive family}

The contrast is the whole point. The three binary nodes each yield $2$ vertices per
parent config in $O(1)$ work. The single $4$-state node $C\text{-}D$ yields up to
$11$ vertices for \emph{one} parent configuration --- and it has $2$ parent configs, so
$4+11 = 15$ vertices come from this one family alone. The Upper Bound Theorem
(next section) caps a $d'{=}3$, $m{=}8$ polytope at $12$ vertices; the observed $11$ is
right at that ceiling. Fusing $C$ and $D$ into a cardinality-$4$ node --- forced by the
cross-atom sentence $s_4$ --- is what pushed the enumeration from trivial to non-trivial.

\section{Complexity for larger credal networks}

\subsection{Two sources of blow-up}

Equation~\eqref{eq:lrs-cost} is polynomial \emph{in the output}, so the real question is
how large the output $V$ is, and how many polytopes must be enumerated. There are two
independent exponential factors, both driven by the \textbf{family scope}.

\paragraph{(1) Vertices per polytope, exponential in cardinality $d$.}
By the \textbf{Upper Bound Theorem} (McMullen), a $d'$-dimensional polytope with $m$
facets has at most
\[
  V_{\max}(d', m) \;=\; \binom{m - \lceil d'/2\rceil}{m - d'}
    \;+\; \binom{m - \lceil (d'{+}1)/2\rceil}{m - d'}
\]
vertices. For a local interval simplex, $d' = d-1$ and $m = 2d$, so $V$ grows
\emph{super-polynomially in the node cardinality $d$}:
\begin{center}
\begin{tabular}{rrrr}
\toprule
states $d$ & polytope dim $d'$ & facets $m$ & max vertices $V_{\max}$ \\
\midrule
2  & 1  & 4  & 2 \\
4  & 3  & 8  & 12 \\
8  & 7  & 16 & 440 \\
16 & 15 & 32 & 692{,}208 \\
\bottomrule
\end{tabular}
\end{center}
A node of cardinality $d = 2^{k}$ arises whenever $k$ atoms are fused into one compound
node (as $C,D \to C\text{-}D$ with $k=2$, $d=4$ above). Thus \emph{each additional fused
atom roughly squares the worst-case vertex count} of that family's local polytope.

\paragraph{(2) Number of polytopes, exponential in parent scope.}
LRS is invoked once per parent configuration. A family whose parents have $p$ atoms in
total has $2^{p}$ parent configurations, hence $2^{p}$ separate polytopes to enumerate.
The total enumeration cost of the whole network is therefore
\begin{equation}
  \underbrace{\sum_{\text{node } X}}_{\text{linear in \#families}}\
  \underbrace{\sum_{\pi \in \mathrm{cfg}(\mathrm{Pa}(X))}}_{\;2^{\,|\mathrm{Pa\ atoms}|}\;}
  \ O\bigl(d'_X \cdot m_X \cdot V_{X,\pi}\bigr),
  \label{eq:total}
\end{equation}
which is \textbf{linear in the number of families} but \textbf{exponential in each
family's scope} through both the $2^{p}$ configurations and the $V_{X,\pi}$ vertices per
configuration (itself up to $\approx V_{\max}(2^{c}{-}1, 2^{c+1})$ for a $c$-atom child).

\subsection{The controlling parameter is treewidth, not $n$}

Because the exponents in~\eqref{eq:total} are the child cardinality and the parent-scope,
and both are bounded by the size of the largest family, LRS enumeration cost is governed
by the \textbf{induced width / treewidth} of the moralized network, \emph{not} by the raw
number of variables $n$:
\begin{itemize}
  \item \textbf{Chains, trees, polytrees} (small, fixed family scope): every family has
    $O(1)$ atoms, so each polytope is low-dimensional with $O(1)$ vertices, and the total
    cost is $O(n)$ --- enumeration is essentially free.
  \item \textbf{$k$-trees} (bounded treewidth $k$): each family conditions on a $k$-clique,
    so a family spans $O(k)$ atoms, giving $O(2^{k})$ parent configs and up to
    $V_{\max}(O(2^{k}))$ vertices each. Cost is \emph{linear in $n$ but exponential in
    $k$} --- exactly why the \texttt{ktree}/\texttt{ktree-fr} benchmark generator
    parametrizes by treewidth $k$.
  \item \textbf{Dense / heavily-merged networks} (large families, e.g. under a large
    \texttt{merge\_budget}, scheme D2): a merged super-family of flattened scope $w$ has a
    child of cardinality up to $2^{w}$ and $2^{w}$-ish parent configs. Both exponentials
    fire simultaneously and the enumeration dominates the entire pipeline --- doubly
    exponential in $w$ in the worst case.
\end{itemize}

\subsection{Practical consequences in this codebase}

\begin{itemize}
  \item \textbf{The \texttt{enumerate\_vertices=False} escape hatch.} CredalJT (scheme D5,
    \texttt{junction\_nlp.py}) formulates its NLP directly from the interval bounds and the
    LMC equalities, so it \emph{skips} LRS entirely (see
    \texttt{CredalNetworkVertices.\_build}). For wide families this avoids the
    exponential~\eqref{eq:total} altogether --- the intervals are $O(2^{\text{scope}})$ to
    store but the vertex set they would generate is far larger.
  \item \textbf{Potential clustering.} CredalVE/CredalCTE offer $n\_clusters$ /
    \texttt{cluster\_representative} knobs that \emph{prune} the enumerated vertex set to a
    bounded number of representatives, trading tightness for a hard cap on the $V$ that
    propagates downstream --- a direct mitigation of factor~(1).
  \item \textbf{The \texttt{.cn} cache pays this cost on every hit.} Loading a compiled
    \texttt{.cn} restores the interval bounds for free, but the engine still re-runs LRS
    (the intervals, not the vertices, are cached). Hence a cache hit's reported build time
    is the stored \texttt{compile\_time} \emph{plus} local vertex enumeration --- cheap for
    narrow families, non-trivial for wide ones. Caching the V-representation itself is
    possible but was rejected: it can be exponentially larger than the H-representation and
    would couple the \texttt{.cn} to pyAgrum.
\end{itemize}

\section{Summary}

LRS is an \emph{output-sensitive, polynomial-per-vertex, constant-space} vertex-enumeration
algorithm: enumerating a local credal set with $V$ vertices costs $O(d' m V)$ time and
$O(d' m)$ memory (Eq.~\eqref{eq:lrs-cost}). The algorithm itself is not the bottleneck ---
the \emph{number of vertices} $V$ and the \emph{number of local credal sets} are. Both grow
exponentially in the \textbf{family scope} (Eq.~\eqref{eq:total}): the child cardinality
inflates $V$ super-polynomially (Upper Bound Theorem: $2\to2$, $4\to12$, $8\to440$,
$16\to\sim\!700{,}000$ vertices), and the parent scope multiplies the number of polytopes by
$2^{p}$. Consequently LRS enumeration is linear in $n$ but exponential in treewidth: free on
chains and trees, tolerable on bounded-treewidth $k$-trees, and prohibitive on dense or
heavily merged networks --- for which the intended remedy is to skip enumeration (CredalJT)
or to cap the vertex set (potential clustering).

\begin{thebibliography}{9}
\bibitem{avisfukuda}
  D.~Avis and K.~Fukuda.
  \emph{A pivoting algorithm for convex hulls and vertex enumeration of arrangements
  and polyhedra.}
  Discrete \& Computational Geometry, 8:295--313, 1992.
\bibitem{avis2000}
  D.~Avis.
  \emph{A revised implementation of the reverse search vertex enumeration algorithm.}
  In Polytopes---Combinatorics and Computation, pp.~177--198, Birkh\"auser, 2000.
\end{thebibliography}

\end{document}
